\begingroup
\setstretch{1}
\begin{tikzpicture}[
  font=\sffamily\small,
  >=Stealth,
  block/.style={draw=colEntorno!80!black,fill=colEntorno!13,rounded corners=3pt,
    text width=3.25cm,minimum height=1.05cm,align=center,inner sep=5pt},
  interface/.style={block,fill=colInterfaz!20,draw=colInterfaz!75!black},
  learn/.style={block,fill=colAgente!16,draw=colAgente!75!black},
  persist/.style={block,fill=catLilaSuave!17,draw=catLilaSuave!80!black},
  flow/.style={-{Stealth[length=2.1mm]},line width=.8pt,draw=colEntorno!80!black},
  return/.style={-{Stealth[length=2.1mm]},line width=.8pt,draw=colAgente!80!black}
]
\node[block] (sumo) at (0,0)
  {SUMO / libsumo};
\node[interface] (env) at (0,-2.3)
  {\texttt{SumoRllibEnv}};
\node[block] (workers) at (6.5,-2.3)
  {Muestreo (RLlib)};
\node[learn] (learner) at (6.5,-4.8)
  {Aprendizaje (PyTorch)};
\node[persist] (archive) at (0,-4.8)
  {Almacenamiento};

% Flujos de ida y vuelta separados, con entradas perpendiculares a los bloques.
\draw[flow] ($(sumo.south)+(.9,0)$) --
  node[right,font=\scriptsize,inner sep=2pt] {mediciones}
  ($(env.north)+(.9,0)$);
\draw[return] ($(env.north)+(-.9,0)$) --
  node[left,font=\scriptsize,align=right,inner sep=3pt] {fases y\\avance}
  ($(sumo.south)+(-.9,0)$);
\draw[flow] ($(env.east)+(0,.4)$) --
  node[above,font=\scriptsize,align=center,inner sep=3pt]
  {datos}
  ($(workers.west)+(0,.4)$);
\draw[return] ($(workers.west)+(0,-.4)$) --
  node[below,font=\scriptsize,inner sep=3pt] {acciones}
  ($(env.east)+(0,-.4)$);
\draw[flow] ($(workers.south)+(-.9,0)$) --
  node[left,font=\scriptsize,inner sep=3pt] {transiciones}
  ($(learner.north)+(-.9,0)$);
\draw[return] ($(learner.north)+(.9,0)$) --
  node[right,font=\scriptsize,inner sep=3pt] {parámetros}
  ($(workers.south)+(.9,0)$);
% El recorrido exterior evita atravesar el entorno y entra por archive.west.
\coordinate (storageway) at (-2.8,-4.8);
\draw[flow] (sumo.west) -- (-2.8,0) --
  node[pos=.72,right,font=\scriptsize,inner sep=3pt] {estado}
  (storageway) -- (archive.west);
\end{tikzpicture}
\endgroup
